\begin{vidu}
Một mạch dao động điện từ LC gồm một tụ điện có giá trị $C=23,6\ pF$ và cuộn cảm có độ tự cảm $L=1\ mF$. Trong không khí, tốc độ truyền sóng điện từ là $3.10^8\ m/s$.
\begin{itemize}
\item[a)] Tính bước sóng mà mạch trên thu được.
\item[b)] Sóng điện từ do mạch dao động này phát ra thuộc loại sóng nào?
\item[c)] Bố trí mạch dao động này ở lối vào của máy thu thanh và thay tụ điện bằng một tụ khác có điện dung biến thiên trong khoảng từ $1\ pF$ đến $50\ pF$. Máy thu này có thể thu được sóng điện từ có bước sóng trong khoảng nào?
\item[d)] Mạch này phát ra sóng điện từ lan truyền trong chân không dọc theo đường thẳng từ M đến N cách nhau 217,5 m. Ở thời điểm t, cường độ điện trường tại M bằng 0. Thời gian ngắn nhất sau đó bao lâu thì cường độ điện trường tại N bằng 0?
\end{itemize}	
\loigiai{
\begin{itemize}
\item[a)] Bước sóng vô tuyến mà mạch thu được: $\lambda =2\pi c\sqrt{LC}\approx 290\ m$.
\item[b)] $100\ m \geq \lambda \geq 1000\ m$ nên sóng này thuộc loại sóng trung.
\item[c)] Khoảng giá trị bước sóng mà mạch có thể cộng hưởng được là:
\begin{align*}
&2\pi c\sqrt{L{C_{\min}}}\le \lambda \le 2\pi c\sqrt{LC_{max}}\\
\Leftrightarrow\ &2\pi {.3.10^8}\sqrt{{10^{-3}}{.10.10^{-12}}}\le \lambda \le 2\pi {.3.10^8}\sqrt{{10^{-3}}{.50.10^{-12}}}
\Leftrightarrow 59,6\ m\le \lambda \le 421,5\ m.
\end{align*}
\item[d)] $\Delta \varphi _{M/N}=\dfrac{2\pi MN}{\lambda }=\dfrac{3\pi }{2}\Rightarrow$ M nhanh pha hơn N một góc $\dfrac{3\pi }{2}$.\\
Tại: $t\Rightarrow E_M=0\Rightarrow E_M=\pm E_0\Rightarrow$ thời điểm gần nhất $E_N=0$ là $t=\dfrac{T}{4}=2,4.10^{-7}\ s$.
\end{itemize}
}
\end{vidu}

\loai{Tính bước sóng khi biết $I_0$ và $Q_0$}
\begin{vidu} %[3P4-4.2B][Loai2]  
Mạch chọn sóng của máy thu vô tuyến điện đang hoạt động, người ta đo được $I_0=10\ A$ và $Q_0=10^{-5}\ C$.  Mạch đang dao động bắt được sóng có bước sóng
\choice
{188,5 m}
{18,85 m}
{18,85 km}
{\True 1885 m}
\loigiai{
\begin{itemize}
\item Theo đề bài ta tính được tần số góc: $\omega =\dfrac{I_0}{Q_0}=10^6\ rad/s$.
\item Mạch dao động sẽ bắt được sóng có bước sóng: $\lambda =c.T=c.\dfrac{2\pi}{\omega}=3.10^8.\dfrac{2\pi}{10^6}\approx 1885\ m$.
\end{itemize}
}
\end{vidu}
\loai{Tính bước sóng khi biết T, f}
\begin{vidu} %[3P4-4.2Y][Loai3]
\nguon{QG - 2020} Một sóng điện từ có tần số 120 kHz đang lan truyền trong chân không. Lấy $c = 3.10^8\ m/s$). Sóng này có bước sóng là
\choice
{\True 2500\ m}
{0,8\ m}
{1250\ m}	
{0,4\ m}
\loigiai{
$\lambda =\dfrac{v}{f}=\dfrac{3.10^8}{120.10^3}= 2500\ m$}
\end{vidu}
\loai{Tính dải bước sóng}
\begin{vidu} %[3P4-4.2B][Loai7]
\nguon{QG - 2017} Mạch dao động ở lối vào của một máy thu thanh gồm cuộn cảm thuần có độ tự cảm $3\ \mu H$ và tụ điện có điện dung biến thiên trong khoảng từ 10 pF đến 500 pF. Biết rằng, muốn thu được sóng điện từ thì tần số riêng của mạch phải bằng với tần số của sóng điện từ cần thu (để có cộng hưởng). Trong không khí, tốc độ truyền sóng điện từ là $3.10^8$ m/s, máy thu này có thể thu được các sóng điện từ có bước sóng nằm trong khoảng
\choice
{10 m đến 730 m}
{100 m đến 730 m}
{\True 10 m đến 73 m}
{1 m đến 73 m}
\loigiai{
Khoảng giá trị bước sóng mà mạch có thể cộng hưởng được là:
$2\pi c\sqrt{L{C_{\min}}}\le \lambda \le 2\pi c\sqrt{LC_{max}}$\\
$\Leftrightarrow 2\pi {.3.10^8}\sqrt{{3.10^{-6}}{.10.10^{-12}}}\le \lambda \le 2\pi {.3.10^8}\sqrt{{3.10^{-6}}{.500.10^{-12}}}
\Leftrightarrow 10\ m\le \lambda \le 73\ m$.
}
\end{vidu}
\loai{Xác định bước sóng thuộc vùng sóng nào}
\begin{vidu} %[3P4-4.2B][Loai6] 
Một mạch dao động điện từ LC có điện tích cực đại trên bản tụ là \bth{1\ \mu C} và dòng điện cực đại qua cuộn dây là 0,314 A. Sóng điện từ  do mạch dao động này phát ra thuộc loại
\choice
{\True sóng dài và cực dài}
{sóng trung}
{sóng ngắn}
{sóng cực ngắn}
\loigiai{
\begin{itemize}
\item Ta có: \bth{I_0=\omega Q_0=2\pi fQ_0\Rightarrow f=\dfrac{I_0}{2\pi Q_0}\ct{ và } \lambda=cT=\dfrac{c}{f}}.
\item Suy ra bước sóng điện từ
\begin{align*}
\lambda=c\dfrac{2\pi Q_0}{I_0}=3.10^{8}\dfrac{2.3,14.1.10^{-6}}{0,314}=6003\ m.
\end{align*}
\item Giá trị của \bth{\lambda} thuộc vùng sóng dài và cực dài.
\end{itemize}
}
\end{vidu}
\loai{Tính T, f}
\begin{vidu} %[3P4-4.2B][Loai4]
\nguon{QG - 2019} Một sóng điện từ lan truyền trong chân không có bước sóng 3000 m. Lấy $c=3.10^8$ m/s. Biết trong sóng điện từ, thành phần từ trường tại một điểm biến thiên điều hòa với chu kì T. Giá trị của T là
\choice
{$4.10^{-6}$ s}
{$2.10^{-5}$ s}
{\True $10^{-5}$ s}
{$3.10^{-6}$ s}
\loigiai{Chu kì của sóng: $T=\dfrac{\lambda}{c}=10^{-5}\ s$.
}
\end{vidu}
\loai{Tính L, C}
\begin{vidu} %[3P4-4.2K][Loai5]
Mạch chọn sóng của một máy thu gồm một tụ điện có điện dung $C = \dfrac{4}{9\pi^2}$ pF và cuộn cảm có độ tự cảm biến thiên. Để có thể bắt được sóng điện từ có bước sóng $\lambda = 100$ m thì độ tự cảm cuộn dây bằng
\choice
{$L = 0,0645$ H}
{\True $L = 0,0625$ H}
{$L = 0,0615$ H}
{$L = 0,0635$ H}
\loigiai{
Độ tự cảm của cuộn dây là: $L=\dfrac{{\lambda}^2}{4{\pi}^2c^2C}=0,0625\ H $. 
}
\end{vidu}
\loai{Tính dải C hoặc L khi biết dải bước sóng}
\begin{vidu} %[3P4-4.2B][Loai8]  
\nguon{QG - 2017} Mạch dao động ở lối vào của một máy thu thanh gồm cuộn cảm thuần có độ tự cảm 5 $\mu H$ và tụ điện có điện dung thay đổi được. Biết rằng, muốn thu được sóng điện từ thì tần số riêng của mạch dao động phải bằng tần số của sóng điện từ cần thu (để có cộng hưởng). Trong không khí, tốc độ truyền sóng điện từ là $3.10^8$ m/s, để thu được sóng điện từ có bước sóng từ 40 m đến 1000 m thì phải điều chỉnh điện dung của tụ điện có giá trị
\choice
{\True từ 90 pF đến 56,3 nF}
{từ 9 pF đến 56,3 nF}
{từ 90 pF đến 5,63 nF}
{từ 9 pF đến 5,63 nF}
\loigiai{
Để thu được sóng thì phải có cộng hưởng $f_0=f\Leftrightarrow \dfrac{1}{2\pi \sqrt{LC}}=\dfrac{c}{\lambda}\Rightarrow \lambda =c.2\pi \sqrt{LC}\Rightarrow C=\dfrac{\lambda^2}{\left(c.2\pi  \right)^2L}$\\
$\Rightarrow \left\{\begin{aligned}
& {C_{\min}}=\dfrac{{40^2}}{{\left({3.10^8}.2\pi  \right)^2}{.5.10^{-6}}}={9.10^{-11}}\ F =90\ pF \\ 
& {C_{\max}}=\dfrac{{100^2}}{{\left({3.10^8}.2\pi  \right)^2}{.5.10^{-6}}}=5,{63.10^{-8}}\ F=56,3\ nF \\ 
\end{aligned} \right.$
}
\end{vidu}
\loai{Trộn sóng trước khi truyền tải}
\begin{vidu} %[3P4-4.2K][Loai9]  
\nguon{ĐH - 2010} Trong thông tin liên lạc bằng sóng vô tuyến, người ta sử dụng cách biến điệu biên độ, tức là làm cho biên độ của sóng điện từ cao tần (gọi là sóng mang) biến thiên theo thời gian với tần số bằng tần số của dao động âm tần. Cho tần số sóng mang là 800 kHz. Khi dao động âm tần có tần số 1000 Hz thực hiện một dao động toàn phần thì dao động cao tần thực hiện được số dao động toàn phần là
\choice
{\True 800}
{1000}
{625}
{1600}
\loigiai{
\begin{itemize}
\item Tần số sóng mang chính là tần số của sóng cao tần.
\item Số dao động thực hiện được trong khoảng thời gian $\Delta t$: $N=f.\Delta t$.
\item Xét trong cùng khoảng thời gian $\Delta t$ thì: $\dfrac{N_c}{N_a}=\dfrac{f_c}{f_a}\Leftrightarrow N_c=N_a.\dfrac{f_c}{f_a}=1.\dfrac{800000}{1000}=800$ dao  động.
\end{itemize}
}
\end{vidu}
\loai{Bài toán liên quan tụ điện phẳng}
\begin{vidu} %[3P4-4.2K][Loai10]
Mạch thu sóng có $L = 1$ mH, tụ điện C là tụ phẳng không khí có diện tích đối diện $40\ cm^2$, khoảng cách giữa hai bản tụ là 1,5 mm. Bước sóng mà mạch thu được là
\choice
{\True 289 m}
{354 m}
{298 m}
{453 m}
\loigiai{
Ta có: $C=\dfrac{\varepsilon S}{4kd\pi}=2,36.10^{-11}\ F\Rightarrow \lambda =2\pi c\sqrt{LC}=289\ m$.
}
\end{vidu}
\loai{Thay đổi C}
\begin{vidu} %[3P4-4.2B][Loai11]
Một mạch chọn sóng gồm cuộn dây có hệ số tự cảm không đổi và một tụ điện có điện dung biến thiên. Khi điện dung của tụ là 50 nF thì mạch thu được bước sóng $\lambda  = 50$ m. Nếu muốn thu được bước sóng $\lambda  = 30$ m thì phải điều chỉnh điện dung của tụ
\choice
{giảm 30 nF}
{\True giảm 32 nF}
{giảm 25 nF}
{giảm 18 nF}
\loigiai{
\begin{itemize}
\item Ta có: $\lambda =c2\pi \sqrt{LC}\Rightarrow C\sim \lambda^2$.
\item Có $\lambda_2=\dfrac{30}{50}{\lambda_1}=0,6{\lambda_1}\Rightarrow {C_2}={{(0,6)}^2}{C_1}={{(0,6)}^2}.50=18\ nF $.
\item Điện dung của tụ giảm $50-18=32$ nF. 
\end{itemize}
}
\end{vidu}
\loai{Thay đổi cả L và C}
\begin{vidu} %[3P4-4.2K][Loai12]
Mạch chọn sóng của một máy thu vô tuyến điện gồm bộ tụ điện và cuộn cảm thuần L. Khi $L = L_1;\ C = C_1$ thì mạch thu được bước sóng $\lambda$. Khi $L = 3L_1;\ C = C_2$ thì mạch thu được bước sóng là $2\lambda$. Khi điều chỉnh cho $L = 3L_1;\ C = C_1 + 2C_2$ thì mạch thu được bước sóng là
\choice
{$\lambda\sqrt{10}$}
{\True $\lambda\sqrt{11}$}
{$\lambda\sqrt{5}$}
{$\lambda\sqrt{7}$}
\loigiai{
\begin{itemize}
\item
Ta có: $\left(\dfrac{{\lambda_1}}{{\lambda_2}} \right)^2=\dfrac{{C_1}{L_1}}{{C_2}{L_2}}\Leftrightarrow {\left(\dfrac{1}{2} \right)^2}=\dfrac{{C_1}}{3{C_2}}\Leftrightarrow \dfrac{{C_1}}{{C_2}}=\dfrac{3}{4}$.
\item 	Mặt khác: $\left(\dfrac{{\lambda_1}}{{\lambda_0}} \right)^2=\dfrac{{L_1}{C_1}}{3{L_1}({C_1}+2{C_2})}\Leftrightarrow {\left(\dfrac{\lambda}{{\lambda_0}} \right)^2}=\dfrac{1}{3}.\dfrac{3}{11}\Leftrightarrow {\lambda_0}=\sqrt{11}\lambda $. 
\end{itemize}
}
\end{vidu}
\loai{Tổng hợp}
\begin{vidu} %[3P4-4.2K][Loai14] 
\nguon{MH - 2019} Một sóng điện từ lan truyền trong chân không dọc theo đường thẳng từ điểm M đến điểm N cách nhau 45 m. Biết sóng này có thành phần điện trường tại mỗi điểm biến thiên điều hòa theo thời gian với tần số 5 MHz. Lấy $c = 3.10^8$ m/s. Ở thời điểm t, cường độ điện trường tại M bằng 0. Thời điểm nào sau đây cường độ điện trường tại N bằng 0?
\choice
{$t + 225$ ns}
{$t + 230$ ns}
{$t + 260$ ns}
{\True $t + 250$ ns}
\loigiai{
\begin{itemize}
\item $\Delta \varphi _{M/N}=\dfrac{2\pi MN}{\lambda }=\dfrac{3\pi }{2}\Rightarrow$ M nhanh pha hơn N một góc $\dfrac{3\pi }{2}$.
\item Tại: $t\Rightarrow E_{M}=0\Rightarrow E_{N}=\pm E_{0}\Rightarrow$ thời điểm $E_N=0$ là $t'=t+m\dfrac{T}{4}=t+50m\ ns$.
\item Với $m = 5$ sẽ thỏa mãn $t'=t+250\ ns$.
\end{itemize}
}
\end{vidu}
\begin{vidu}%Câu 40%E2
Một anten parabol đặt tại điểm O trên mặt đất, phát ra một sóng truyền theo phương làm với mặt phẳng ngang một góc $45^\circ$ hướng lên cao. Sóng này phản xạ trên tầng điện li, rồi trở lại gặp mặt đất ở điểm M. Cho bán kính Trái Đất $R=6400\ km$. Tầng điện li coi như một lớp cầu ở độ cao 100 km trên mặt đất. Cho 1 phút$=3.10^{-4}\ rad$. Độ dài cung OM là
\choice
{\True 201,6 km}
{206,1 km}
{126,0 km}
{162,1 km}
\loigiai{

%\begin{center} \begin{figure}[h]
%	\centering
%	\includegraphics[width=0.3\linewidth]{"../LOC CAU HAY 2021/Hinh/40"}
%\end{figure} \end{center}
\begin{itemize}
\item Để tính độ dài cung AM ta tính góc $\varphi=\angle O{O'}M$\\
\item Xét $\Delta OO’A$: $ O{O'}=R;{O'}A=R+h;\beta=\angle O{O'}A=135^\circ$\\
\item Theo định lý hàm số sin:\\
$\dfrac{O'A}{\sin{135^\circ}}=\dfrac{O{O'}}{\sin\dfrac{\alpha}{2}}\Rightarrow\sin\dfrac{\alpha}{2}=\dfrac{O{O'}}{O'A}\sin{135^\circ}=0,696$\\
$\Rightarrow\alpha=88,25^\circ\Rightarrow\varphi=360^\circ-270^\circ-88,25^\circ$ $=1,75^\circ=1,75.60.3.10^{-4}=315.10^{-4}rad$\\
\item Cung $ AM=R\varphi=315.10^{-4}.6,4.10^3\left(km\right)=201,6km$
\end{itemize}
}
\end{vidu}